\chapter{Experience Paths and Character Building}
łabel{ch:xp−paths}
∈dex{character building}∈dex{XP!paths}

How you spend your \textbf{Experience Points (XP)} shapes more than your
numbers on a sheet−−−it defines your character's \emph{career}, their growth
over time, and the kind of stories they invite into play.

This chapter focuses on:
\begin{itemize}
  ℹtem turning a loose idea into a clear \textbf{character concept},
  ℹtem mapping that concept into \textbf{starting choices} that fit the campaign rules,
  ℹtem choosing a \textbf{career path} (how you plan to grow),
  ℹtem and understanding how your character can evolve across \textbf{early, mid, and late} game.
\end{itemize}

All examples in this chapter are \emph{legal starting builds} under the standard
creation rules, and are written to be directly usable by players.

% ============================================================
% FROM CONCEPT TO CAREER
% ============================================================

\section{From Concept to Career}
łabel{sec:xp−concept−to−career}
∈dex{concept mapping}∈dex{career planning}

Before you spend a single XP, start with the story you want to tell.

\subsection*{Step 1: Concept Seed}
łabel{sec:xp−concept−seed}

Begin with a simple statement of who your character is:
\begin{itemize}
  ℹtem ``Stoic bodyguard who will take the hit for anyone else.''
  ℹtem ``Nervy courier who knows every back alley in the city.''
  ℹtem ``Soft−spoken merchant who pulls strings behind the scenes.''
\end{itemize}

This is your \textbf{concept seed}. Everything else grows from here.

\subsection*{Step 2: Role in the Party}
łabel{sec:xp−role−party}

Decide what you want to \emph{do at the table}:
\begin{itemize}
  ℹtem Take hits and lock down threats?
  ℹtem Spot danger and solve practical problems?
  ℹtem Talk through conflicts and manage social fallout?
  ℹtem Create tools, magic, or plans that enable everyone else?
\end{itemize}

Your role helps you choose a \textbf{Primary Attribute} and 2−−3 \textbf{Key Skills}
(see \ref{ch:advancement}).

\subsection*{Step 3: Career Theme}
łabel{sec:xp−career−theme}

Ask yourself:
\begin{itemize}
  ℹtem Is this a story of \emph{self−mastery}? (training, discipline, pushing limits)
  ℹtem A story of \emph{connections}? (alliances, obligations, webs of influence)
  ℹtem A story of \emph{balance}? (learning to juggle many demands at once)
\end{itemize}

Your answer naturally points you toward one of the three \textbf{XP paths} in
this chapter.

\subsection*{Step 4: Long Arc in One Sentence}
łabel{sec:xp−long−arc}

Write a single sentence for where you imagine your character at high XP:
\begin{itemize}
  ℹtem ``I want to become the greatest duelist in the realm.''
  ℹtem ``I want to map the unknown and bring others safely home.''
  ℹtem ``I want to run the network everyone relies on but no one sees.''
\end{itemize}

You do not have to stick to this forever, but it gives you a \emph{direction}
for your early choices.

\begin{tcolorbox}[colback=blue!5!white,colframe=blue!75!black,title=Character Career Checklist,fonttitle=\bfseries]
\textbf{Before spending XP, decide:}
\begin{enumerate}
  ℹtem Your \textbf{concept seed} (one sentence).
  ℹtem Your \textbf{role in the party} (what you do most).
  ℹtem Your preferred \textbf{career theme}: self, balance, or influence.
  ℹtem A rough \textbf{long arc} (where you want to end up).
\end{enumerate}
Start simple. You can always refine these as the campaign unfolds.
\end{tcolorbox}

% ============================================================
% MAPPING CONCEPT TO MECHANICS
% ============================================================

\section{Mapping Concept to Mechanics}
łabel{sec:xp−mapping}
∈dex{concept mapping!to mechanics}

Once you have a concept, you turn it into numbers. Use this as a guideline:

\begin{enumerate}
  ℹtem Pick a \textbf{Primary Attribute} that supports your role
    (Body, Wits, Spirit, or Presence) −− see \ref{sec:core−attributes}.
  ℹtem Choose 2−−3 \textbf{Key Skills} that you want to be noticeably good at
    (see \ref{sec:skills}).
  ℹtem Choose a \textbf{Career Path} (Personal, Balanced, Influencer) that matches
        your theme.
  ℹtem Use the starting XP budget in \ref{sec:xp−starting−guidelines} to buy:
    \begin{itemize}
      ℹtem Attribute steps for your primary and possibly secondary stats.
      ℹtem Skill steps for your key and supporting skills.
      ℹtem Optional resources or special abilities, if your path calls for them.
    \end{itemize}
\end{enumerate}

Throughout this chapter, each example shows how a concept becomes a legal build.

% ============================================================
% THREE CAREER PATHS
% ============================================================

\section{Three Career Paths}
łabel{sec:xp−three−paths}
∈dex{advancement paths}∈dex{career paths}

There are three broad approaches to character development, each representing a
different \textbf{career philosophy}:

\begin{description}
ℹtem[Personal Path] Focus on personal mastery and self−improvement (\ref{sec:xp−personal−path}).
ℹtem[Balanced Path] Mix personal growth with resources and influence (\ref{sec:xp−balanced−path}).
ℹtem[Influencer Path] Build networks, assets, and strategic power (\ref{sec:xp−influencer−path}).
\end{description}

These paths are not rigid classes. They are \emph{patterns} that help you
decide where most of your XP will go, especially in the early game.

% ============================================================
% PLAYER ARCHETYPES
% ============================================================

\section{Player Archetypes}
łabel{sec:player−archetypes}
∈dex{player archetypes}∈dex{Solo}∈dex{Mixed}∈dex{Mastermind}

Not every player spends XP the same way. These three \textbf{player archetypes}
describe common patterns. Identify yours −− it will help you choose a career
path and avoid frustration.

\subsection*{The Solo}
łabel{sec:archetype−solo}

\textbf{Definition:} Invests XP primarily in Attributes and Skills. Minimal
followers, minimal holdings. All power is on the character sheet.

\medskip
\noindent\textbf{Typical XP Spread:} 70−−90% Self; 0−−10% On−screen help; 0−−20% Off−screen assets.

\medskip
\noindent\textbf{Strengths:} Consistent scene impact; few moving parts; resilient to
follower loss.

\medskip
\noindent\textbf{Risks:} Limited fiction reach between sessions; can stall when
problems demand logistics or networks.

\medskip
\noindent\textbf{Starting Focus:} Attributes 2−−3, Skills 1−−2, minimal assets.

\subsection*{The Mixed Player}
łabel{sec:archetype−mixed}

\textbf{Definition:} Splits XP between self−growth and one or two meaningful
assets (a small follower or a reliable holding).

\medskip
\noindent\textbf{Typical XP Spread:} 50−−65% Self; 15−−░ On−screen help; 15−−░ Off−screen assets.

\medskip
\noindent\textbf{Strengths:} Versatile: credible in scenes and has a lever for special
problems.

\medskip
\noindent\textbf{Risks:} Upkeep pressure; helper can be targeted when the GM spends 2+ Story Beats.

\medskip
\noindent\textbf{Starting Focus:} Balanced approach with one minor asset or low−cap follower.

\subsection*{The Mastermind}
łabel{sec:archetype−mastermind}

\textbf{Definition:} Prioritises followers / cadres / familiars and off−screen
networks. The character sheet is the hub of a larger apparatus.

\medskip
\noindent\textbf{Typical XP Spread:} 25−−40% Self; 35−−55% On−screen help; 20−−40% Off−screen assets.

\medskip
\noindent\textbf{Strengths:} Scene control via assistance; strategic reach between
sessions; strong heist / social−planning play.

\medskip
\noindent\textbf{Risks:} Dependency on assist lanes; followers can be endangered on 2+ SB
spends; upkeep pressure.

\medskip
\noindent\textbf{Starting Focus:} Significant investment in followers or major assets
from the beginning.

\begin{tcolorbox}[colback=cyan!5!white,colframe=cyan!75!black,title=Which Archetype Are You?,fonttitle=\bfseries]
Ask yourself:
\begin{itemize}
  ℹtem Do I want my character to be self−sufficient and always ready? → \textbf{Solo}
  ℹtem Do I enjoy having a helpful sidekick or a home base, but not too much bookkeeping? → \textbf{Mixed}
  ℹtem Do I love managing networks, commanding troops, and solving problems through influence? → \textbf{Mastermind}
\end{itemize}
There is no wrong answer −− but knowing your style helps you build a character you will enjoy for many sessions.
\end{tcolorbox}

% ============================================================
% PATH 1: PERSONAL DEVELOPMENT
% ============================================================

\section{Path 1: Personal Development}
łabel{sec:xp−personal−path}
∈dex{personal path}

The \textbf{Personal Path} centers your character's career on what \emph{they}
can do: their body, mind, and honed skills. This is the path of the duelist,
the ascetic, the elite scout, the master healer.

\subsection*{Typical Investment}
\begin{itemize}
ℹtem 70−−90% Personal improvement (Attributes and Skills)
ℹtem 0−−10% Resources and assets
ℹtem 0−−20% Special abilities
\end{itemize}

\subsection*{Strengths}
\begin{itemize}
ℹtem Reliable in direct challenges and combat
ℹtem Minimal upkeep or management required
ℹtem Resilient to loss of external resources
ℹtem Consistent performance in spotlight moments
\end{itemize}

\subsection*{Weaknesses}
\begin{itemize}
ℹtem Limited influence in social or strategic scenes
ℹtem May struggle with problems requiring networks
ℹtem Less capable in logistics or large−scale operations
ℹtem Dependent on personal presence for all solutions
\end{itemize}

\subsection*{Career View}

On this path, your ``career ladder'' is mostly about pushing numbers related to
\textbf{you}: raising Attributes, deepening Skills, and picking a few defining
abilities. Your late−game questions are:
\begin{itemize}
  ℹtem ``What does a legendary version of this look like?''
  ℹtem ``What scenes prove that I have truly mastered my craft?''
\end{itemize}

\subsection*{Build Example: The Duelist (Legal Start)}

\textbf{Concept Seed:} ``A proud blade whose entire life is training and challenge.''

\medskip

\noindent\textbf{Total XP: 30}   (\textbf{34} with +4 from Bonds/Complications; see \ref{sec:xp−starting−guidelines})

\begin{itemize}
ℹtem \textbf{Attributes}∈dex{attributes}: Body 3, Wits 2, Spirit 1, Presence 1
  \begin{itemize}
  ℹtem Costs (Attributes cost \emph{new rating} \(× 3\)):  
        Body 1\(→\)2 (6), 2\(→\)3 (9) = \textbf{15};  
        Wits 1\(→\)2 (6) = \textbf{6};  
        Spirit/Presence remain 1 = \textbf{0}.  
        \emph{Subtotal: 21 XP}
  \end{itemize}
ℹtem \textbf{Skills}∈dex{skills}: Melee 2, Athletics 1
  \begin{itemize}
  ℹtem Costs (Skills cost \emph{new level} \(× 2\)):  
        Melee 0\(→\)1 (2), 1\(→\)2 (4) = \textbf{6};  
        Athletics 0\(→\)1 = \textbf{2}.  
        \emph{Subtotal: 8 XP}
  \end{itemize}
ℹtem \textbf{Totals}: 21 + 8 = \textbf{29 XP}. Bank \textbf{1 XP}.
ℹtem \textbf{With +4 XP (Bonds/Complications)}:  
  Add \emph{Perception 0\(→\)1} (2) and spend banked 1 XP on \emph{Stealth 0\(→\)1} (2),  
  or instead take \emph{Perception 0\(→\)1} (2) and \emph{Sway 0\(→\)1} (2) for broader utility.  
  \emph{Cap: 36XP}.
\end{itemize}

\noindent\textbf{Career Hook:} Early scenes show duels and drills. Mid game focuses on
notorious opponents. Late game tells the story of a blade whose name carries
across borders.

% ============================================================
% PATH 2: BALANCED APPROACH
% ============================================================

\section{Path 2: Balanced Approach}
łabel{sec:xp−balanced−path}
∈dex{balanced path}

The \textbf{Balanced Path} mixes personal capability with supportive resources.
This is the path of the scout−captain, the field medic with a small clinic, the
problem−solver who always seems to have \emph{something} ready.

\subsection*{Typical Investment}
\begin{itemize}
ℹtem 50−−65% Personal improvement
ℹtem 15−−░ Resources and assets
ℹtem 15−−░ Special abilities
\end{itemize}

\subsection*{Strengths}
\begin{itemize}
ℹtem Adaptable to diverse situations
ℹtem Handles both direct and indirect challenges
ℹtem Excellent supporting role for the group
ℹtem Moderate risk profile
\end{itemize}

\subsection*{Weaknesses}
\begin{itemize}
ℹtem Not exceptional in any single area
ℹtem Requires management of resources
ℹtem Moderate upkeep demands
ℹtem Can be outshone by focused specialists
\end{itemize}

\subsection*{Career View}

Balanced characters often have careers that feel like juggling:
\begin{itemize}
  ℹtem they add contacts, gear, or small bases \emph{and}
  ℹtem keep raising the core skills that let them use those tools well.
\end{itemize}

Their late−game questions are:
\begin{itemize}
  ℹtem ``What does my little network or operation look like at full size?''
  ℹtem ``How do my personal skills and resources fit together as a whole?'' 
\end{itemize}

\subsection*{Build Example: The Scout (Legal Start)}

\textbf{Concept Seed:} ``A pathfinder who knows how to get in, get out, and get others home.''

\medskip

\noindent\textbf{Total XP: 30}   (\textbf{34} with +4 from Bonds/Complications)

\begin{itemize}
ℹtem \textbf{Attributes}: Wits 2, Body 2, Spirit 1, Presence 1
  \begin{itemize}
  ℹtem Costs: Wits 1\(→\)2 (6), Body 1\(→\)2 (6) = \textbf{12 XP}
  \end{itemize}
ℹtem \textbf{Skills}: Survival 2, Perception 1, Stealth 1
  \begin{itemize}
  ℹtem Costs: Survival 0\(→\)1 (2), 1\(→\)2 (4) = \textbf{6};  
        Perception 0\(→\)1 \textbf{2};  
        Stealth 0\(→\)1 \textbf{2}.  
        \emph{Subtotal: 10 XP}
  \end{itemize}
ℹtem \textbf{Resources}∈dex{resources}: \emph{Minor equipment cache} (camp gear, maps, signal kit) = \textbf{4 XP}
ℹtem \textbf{Special Abilities}∈dex{special abilities}: \emph{Wilderness Lore} (broad travel benefits) = \textbf{4 XP}
ℹtem \textbf{Totals}: 12 + 10 + 4 + 4 = \textbf{32 XP}.
ℹtem \textbf{With +4 XP}: add \emph{Perception 1\(→\)2} (+4) \emph{or} take a \emph{trained hawk companion} (Minor Resource, 4 XP).
\end{itemize}

\noindent\textbf{Career Hook:} Early stories show small journeys. Mid game explores
dangerous routes and rescue missions. Late game might see the scout running
expeditions, caravans, or an entire mapping guild.

% ============================================================
% PATH 3: INFLUENCER FOCUS
% ============================================================

\section{Path 3: Influencer Focus}
łabel{sec:xp−influencer−path}
∈dex{influencer path}

The \textbf{Influencer Path} emphasizes networks, assets, and strategic power.
This is the path of the merchant−prince, the fixer, the spymaster, the noble
with more favors owed than coins in their purse.

\subsection*{Typical Investment}
\begin{itemize}
ℹtem 25−−40% Personal improvement
ℹtem 35−−55% Resources and assets
ℹtem 20−−40% Special abilities
\end{itemize}

\subsection*{Strengths}
\begin{itemize}
ℹtem Strong strategic and social influence
ℹtem Can solve problems indirectly
ℹtem Excellent at planning and preparation
ℹtem Creates opportunities for the whole group
\end{itemize}

\subsection*{Weaknesses}
\begin{itemize}
ℹtem Personally vulnerable in direct confrontations
ℹtem High maintenance requirements
ℹtem Complications can cascade through networks
ℹtem Dependent on external factors
\end{itemize}

\subsection*{Career View}

Influencer careers are about building and surviving \textbf{webs}:
\begin{itemize}
  ℹtem webs of contacts, debts, trade routes, cells, or clients.
\end{itemize}
Late game questions become:
\begin{itemize}
  ℹtem ``How big can this web reasonably grow within the fiction?''
  ℹtem ``What happens when a strand of my web is cut, or pulled tight?''
\end{itemize}

\subsection*{Build Example: The Merchant (Legal Start)}

\textbf{Concept Seed:} ``A friendly shopkeeper whose real power lies in who owes them favors.''

\medskip

\noindent\textbf{Total XP: 30}   (\textbf{34} with +4 from Bonds/Complications)

\begin{itemize}
ℹtem \textbf{Attributes}: Presence 2, Wits 2, Spirit 1, Body 1
  \begin{itemize}
  ℹtem Costs: Presence 1\(→\)2 (6), Wits 1\(→\)2 (6) = \textbf{12 XP}
  \end{itemize}
ℹtem \textbf{Skills}: Sway 2, Deception 1, Lore 1
  \begin{itemize}
  ℹtem Costs: Sway 0\(→\)1 (2), 1\(→\)2 (4) = \textbf{6};  
        Deception 0\(→\)1 \textbf{2};  
        Lore 0\(→\)1 \textbf{2}.  
        \emph{Subtotal: 10 XP}
  \end{itemize}
ℹtem \textbf{Resources}: \emph{Standard trading office} (staffed storefront, ledgers, storage) = \textbf{8 XP}
ℹtem \textbf{Totals}: 12 + 10 + 8 = \textbf{32 XP}.
ℹtem \textbf{With +4 XP}: add \emph{Negotiation Mastery} (4 XP general ability) \emph{or} expand to a second \emph{Minor merchant route} (4 XP).
\end{itemize}

\noindent\textbf{Career Hook:} Early scenes are about one shop and a few key clients.
Mid game introduces caravans, deals, and rivals. Late game might place the
merchant at the center of regional politics and supply chains.

% ============================================================
% BUILDING FOR GROUP SYNERGY
% ============================================================

\section{Building for Group Synergy}
łabel{sec:xp−synergy}
∈dex{group synergy}

Your character's career path matters most in context with the rest of the
party. A well−constructed party covers key roles without excessive overlap.

\subsection*{Key Roles to Consider}
\begin{description}
ℹtem[Combat] Frontline defence, skirmishing, or ranged fire support.
ℹtem[Social] Negotiation, intimidation, information gathering.
ℹtem[Exploration] Scouting, navigating, surviving hostile environments.
ℹtem[Magic] Free casting, rites, rituals, or summoning.
ℹtem[Assets] Logistics, networks, safehouses, followers.
\end{description}

\subsection*{Complementary Paths}
\begin{itemize}
ℹtem \textbf{Personal path} characters provide reliable combat and specialist solutions.
ℹtem \textbf{Balanced path} characters handle diverse, connective challenges.
ℹtem \textbf{Influencer path} characters create opportunities and resources for everyone.
\end{itemize}
Together, a mixed party can cover nearly every type of problem.

\subsection*{Redundant Paths}
\begin{itemize}
ℹtem Multiple personal path characters may overlap in combat or role.
ℹtem Multiple influencer path characters may compete for spotlight and assets.
ℹtem Consider diversifying within similar paths: different weapon styles, different Patrons, different social approaches.
ℹtem Example: Two combatants −− one a duelist (single−target damage), one a shield−focused guardian (area protection).
\end{itemize}

\subsection*{Balancing Your Table}
If you notice a gap (e.g., no one has investment in social skills or no one has
a safehouse asset), consider adjusting your XP spending or discussing with the
group. A party that talks about roles early will have a smoother campaign.

% ============================================================
% STARTING GUIDELINES
% ============================================================

\section{Starting Character Guidelines}
łabel{sec:xp−starting−guidelines}
∈dex{character creation!starting XP}

\subsection*{Base XP Allocation}
\begin{itemize}
ℹtem \textbf{Standard Starting XP}: \textbf{32} points
ℹtem \textbf{Bonds and Complications}: You may take up to \textbf{two total} from any mix of meaningful \emph{Bonds} (up to 2, +2 XP each) and significant \emph{Complications} (up to 2, +2 XP each), granting maximum \textbf{+4 XP}. (See \ref{ch:backgrounds} for Bonds and Complications.)
ℹtem \textbf{Maximum Starting XP}: \textbf{36} points
ℹtem \textbf{Complication Effect}: Each unresolved starting Complication adds +1 banked SB to early scenes until cleared.
\end{itemize}

\subsection*{Recommended Starting Ranges}

\begin{center}
\small
\begin{longtable}{ll}
⊤rule
\textbf{Category} & \textbf{Recommended XP} \\
∣rule
Primary Attribute & 9−−12 XP (rating 3−−4) \\
Secondary Attributes & 0−−9 XP each (rating 1−−3) \\
Key Skills & 4−−6 XP each (rating 2−−3) \\
Supporting Skills & 2−−4 XP each (rating 1−−2) \\
Resources & 0−−8 XP total \\
Special Abilities & 0−−8 XP total \\
⊥tomrule
\end{longtable}
\end{center}

\paragraph{Cost Reminders:}
\begin{itemize}
ℹtem \textbf{Attributes}: Each step costs \emph{new rating} \(× 3\) XP (e.g., 1\(→\)2 costs 6; 2\(→\)3 costs 9). See \ref{sec:core−attributes}.
ℹtem \textbf{Skills}: Each step costs \emph{new level} \(× 2\) XP (e.g., 0\(→\)1 costs 2; 1\(→\)2 costs 4). See \ref{sec:skills}.
ℹtem \textbf{Resources}: Minor 4 XP; Standard 8 XP; Major 12 XP.
ℹtem \textbf{Special Abilities}: Minor Edge 2 XP; Major Edge 4 XP; Prestige 6+ XP.
\end{itemize}

% ============================================================
% PROGRESSION PLANNING
% ============================================================

\section{Progression Planning}
łabel{sec:xp−progression}
∈dex{progression planning}

Think of your character's career in phases. These XP bands are approximate and
assume you began at 30−−36XP.

\subsection*{Early Game (0−−40 XP)}
łabel{sec:xp−early−game}

Focus on establishing core capabilities:
\begin{itemize}
ℹtem Reach attribute rating 3 in your primary area.
ℹtem Develop 2−−3 key skills to rating 2−−3.
ℹtem Acquire basic resources or one special ability (if your path uses them).
ℹtem Clarify your character's niche in the group.
\end{itemize}

\noindent\textbf{When to take a Follower vs. Asset vs. Talent (Early):}
\begin{itemize}
  ℹtem \textbf{Follower:} If you want on−screen help (a scout, a bodyguard, an apprentice). Start with Cap 1−−2.
  ℹtem \textbf{Asset:} If you want an off−screen advantage (a safehouse, a workshop, a contact network). Start with Minor Asset (4 XP).
  ℹtem \textbf{Talent:} If you want a signature move (combat technique, social edge, magic access). Minor talents (2−−3 XP) are great early.
\end{itemize}

\subsection*{Mid Game (41−−90 XP)}
łabel{sec:xp−mid−game}

Expand and specialize:
\begin{itemize}
ℹtem Increase your primary attribute to 4.
ℹtem Specialize key skills to rating 3−−4.
ℹtem Develop supporting capabilities that fit your concept.
ℹtem Build strategic resources or networks if on the Balanced or Influencer path.
ℹtem Acquire signature special abilities that define your style.
\end{itemize}

\noindent\textbf{When to take a Follower vs. Asset vs. Talent (Mid):}
\begin{itemize}
  ℹtem \textbf{Follower:} Upgrade a loyal contact to Cap 3 (9 XP) −− they become a reliable asset who can act independently.
  ℹtem \textbf{Asset:} Expand your network −− a Standard Asset (8 XP) turns a safehouse into a spy ring or a workshop into a guild connection.
  ℹtem \textbf{Talent:} Major talents (4−−6 XP) define your mid−game identity −− Backstab, Ritual Mastery, Command Presence.
\end{itemize}

\subsection*{Late Game (91−−150 XP)}
łabel{sec:xp−late−game}

Master your chosen path:
\begin{itemize}
ℹtem Achieve peak attributes (rating 4−−5).
ℹtem Master key skills (rating 4−−5).
ℹtem Build substantial influence or unique capabilities.
ℹtem Develop advanced special abilities.
ℹtem Consider legacy projects, organizations, or long−term changes to the setting.
\end{itemize}

\noindent\textbf{When to take a Follower vs. Asset vs. Talent (Late):}
\begin{itemize}
  ℹtem \textbf{Follower:} A Cap 4−−5 follower (16−−25 XP) can become a successor (see \ref{ch:legacy−engine}) −− invest in someone who could carry on your cause.
  ℹtem \textbf{Asset:} Major Assets (12 XP) give you regional influence −− a fortress, a charter, a fleet.
  ℹtem \textbf{Talent:} Prestige (7−−10 XP) and Epic (11+ XP) talents reshape the story. Plan your build around one or two capstones.
\end{itemize}

\begin{tcolorbox}[colback=green!5!white,colframe=green!75!black,title=XP Path Quick Reference,fonttitle=\bfseries]
\textbf{Personal Path (70−−90% self):}
\begin{itemize}
ℹtem Reliable individual performance.
ℹtem Low upkeep, high consistency.
ℹtem Best for combatants and focused specialists.
\end{itemize}

\textbf{Balanced Path (50−−65% self):}
\begin{itemize}
ℹtem Good all−around capability.
ℹtem Moderate risk and upkeep.
ℹtem Flexible support and problem−solving role.
\end{itemize}

\textbf{Influencer Path (25−−40% self):}
\begin{itemize}
ℹtem Strategic power and influence.
ℹtem High upkeep, high reward.
ℹtem Creates opportunities for the entire group.
\end{itemize}

\textbf{Starting XP:} \textbf{30} base  +  up to \textbf{+4} from Bonds/Complications (max start \textbf{34}).
\end{tcolorbox}

% ============================================================
% PATH COMBINATIONS
% ============================================================

\section{Path Combination Strategies}
łabel{sec:xp−combinations}
∈dex{path combinations}

Many players mix elements from different paths as their character's life
unfolds.

\subsection*{Combat Specialist with Resources}
\begin{itemize}
ℹtem Strong personal combat capabilities.
ℹtem Moderate resource investment for support.
ℹtem Good for frontline fighters who need logistical backup.
ℹtem Example: Warrior with a fortified base and loyal troops.
\end{itemize}

\subsection*{Social Character with Personal Skills}
\begin{itemize}
ℹtem Excellent social capabilities.
ℹtem Solid personal skills for self−defense or utility.
ℹtem Good for diplomats who operate independently.
ℹtem Example: Ambassador with combat training and persuasion skills.
\end{itemize}

\subsection*{Technical Expert with Networks}
\begin{itemize}
ℹtem Deep technical or magical expertise.
ℹtem Network of contacts and resources to enable projects.
ℹtem Good for specialists who need supply chains and helpers.
ℹtem Example: Master crafter with supplier network and apprentices.
\end{itemize}

% ============================================================
% RESOURCE MANAGEMENT
% ============================================================

\section{Resource Management}
łabel{sec:xp−resource−mgmt}
∈dex{resource management}

Each path asks for different management styles across your character's career.

\subsection*{Personal Path Management}
\begin{itemize}
ℹtem Minimal upkeep requirements.
ℹtem Focus on equipment, training scenes, and personal growth.
ℹtem Low complexity, high reliability.
\end{itemize}

\subsection*{Balanced Path Management}
\begin{itemize}
ℹtem Moderate upkeep for resources.
ℹtem Relationship maintenance with contacts.
ℹtem Skill development alongside resource management.
ℹtem Balanced time investment between self and assets.
\end{itemize}

\subsection*{Influencer Path Management}
\begin{itemize}
ℹtem Significant upkeep demands.
ℹtem Network maintenance and expansion.
ℹtem Resource allocation and development.
ℹtem Strategic planning and opportunity management.
\end{itemize}

% ============================================================
% RISK ASSESSMENT
% ============================================================

\section{Risk Assessment}
łabel{sec:xp−risk}
∈dex{risk assessment}

Each path carries different career risks.

\subsection*{Personal Path Risks}
\begin{itemize}
ℹtem Over−specialization in one area.
ℹtem Vulnerability to problems outside your specialty.
ℹtem Limited diversification later in the game.
ℹtem May become predictable in approach.
\end{itemize}

\subsection*{Balanced Path Risks}
\begin{itemize}
ℹtem Jack−of−all−trades, master of none.
ℹtem Spread too thin across capabilities.
ℹtem Moderate risks in multiple areas.
ℹtem May lack standout, defining capabilities if unfocused.
\end{itemize}

\subsection*{Influencer Path Risks}
\begin{itemize}
ℹtem Networks are vulnerable to attack and betrayal.
ℹtem High maintenance requirements.
ℹtem Cascade failure potential if a key asset falls.
ℹtem Personal safety concerns when the web is targeted.
\end{itemize}

% ============================================================
% ADAPTING YOUR PATH
% ============================================================

\section{Adapting Your Path}
łabel{sec:xp−adapting}
∈dex{path adaptation}

Your chosen path is a starting point, not a contract. You can shift focus as
your character's story changes.

\subsection*{Early Shift (0−−40 XP)}
\begin{itemize}
ℹtem Easy to change direction.
ℹtem Minimal sunk cost in any single approach.
ℹtem Good time to experiment with different styles.
ℹtem You can respond quickly to group needs or unexpected story turns.
\end{itemize}

\subsection*{Mid Game Shift (41−−90 XP)}
\begin{itemize}
ℹtem Requires more deliberate planning.
ℹtem Some capabilities need to be maintained while others change.
ℹtem You can fill emerging gaps in group capability.
ℹtem Temporary performance dips may occur during the transition.
\end{itemize}

\subsection*{Late Game Shift (91+ XP)}
\begin{itemize}
ℹtem Significant investment in your current path.
ℹtem Major shifts require substantial XP and story justification.
ℹtem Often better to add complementary capabilities than completely pivot.
ℹtem Consider how changes reflect pivotal events in your character's life.
\end{itemize}

% ============================================================
% PRACTICAL BUILDING EXAMPLES
% ============================================================

\section{Practical Building Examples (Narrative Roles, Legal Starts)}
łabel{sec:xp−examples}
∈dex{build examples}

Each of the following examples combines:
\begin{itemize}
  ℹtem a short \textbf{concept},
  ℹtem a \textbf{path choice},
  ℹtem a legal \textbf{starting build},
  ℹtem and a sketch of their \textbf{career arc}.
\end{itemize}

\subsection*{Example 1: The Guardian}
łabel{sec:xp−ex−guardian}

\textbf{Path}: Personal   \textbf{Total: 32 XP}

\medskip

\noindent\textbf{Concept Seed:} ``Quiet wall of muscle who would rather bleed than lose anyone.''

\begin{itemize}
ℹtem \textbf{Attributes}: Body 3 (1\(→\)2: 6, 2\(→\)3: 9), Wits 2 (1\(→\)2: 6) = \textbf{21 XP}
ℹtem \textbf{Skills}: Melee 2 (0\(→\)1: 2, 1\(→\)2: 4 = 6), Athletics 1 (0\(→\)1: 2) = \textbf{8 XP}
ℹtem \textbf{Bank}: \textbf{1 XP}
ℹtem \textbf{Role at table}: Frontline protection and reliable pressure in fights.
ℹtem \textbf{With +4 XP}: add \emph{Combat Reflexes} (2 XP talent) and \emph{Shield Mastery} (4 XP talent) using banked 1 + 4 = 5 XP (GM may allow rounding or an extra minor Complication), or instead buy more supporting skills.
\end{itemize}

\noindent\textbf{Career Arc:} Early on, the Guardian is ``the one who stands in
front.'' Mid game adds tactical awareness and maybe leadership. Late game might
see them training others or becoming the symbol of a cause.

\subsection*{Example 2: The Explorer}
łabel{sec:xp−ex−explorer}

\textbf{Path}: Balanced   \textbf{Total: 32 XP}

\medskip

\noindent\textbf{Concept Seed:} ``Curious wanderer whose maps are more valuable than gold.''

\begin{itemize}
ℹtem \textbf{Attributes}: Wits 2 (6), Body 2 (6) = \textbf{12 XP}
ℹtem \textbf{Skills}: Survival 2 (0\(→\)1: 2, 1\(→\)2: 4 = 6), Perception 1 (2), Stealth 1 (2) = \textbf{10 XP}
ℹtem \textbf{Resources}: Minor mapping kit & route notes = \textbf{4 XP}
ℹtem \textbf{Ability}: Trail Sense = \textbf{4 XP}
ℹtem \textbf{Totals}: \textbf{32 XP}. With +4 XP, raise \emph{Perception 1\(→\)2} (+4) or add a trained beast (Minor Resource, 4 XP).
\end{itemize}

\noindent\textbf{Career Arc:} Early stories focus on survival and discovery. Mid
game might see them guiding trade, armies, or refugees. Late game, they may
literally redraw the map of the setting.

\subsection*{Example 3: The Schemer}
łabel{sec:xp−ex−schemer}

\textbf{Path}: Influencer   \textbf{Total: 32 XP}

\medskip

\noindent\textbf{Concept Seed:} ``The person who always knows someone who knows someone.''

\begin{itemize}
ℹtem \textbf{Attributes}: Presence 2 (6), Wits 2 (6) = \textbf{12 XP}
ℹtem \textbf{Skills}: Sway 2 (0\(→\)1: 2, 1\(→\)2: 4 = 6), Deception 1 (2), Lore 1 (2) = \textbf{10 XP}
ℹtem \textbf{Resources}: Standard safehouse & message drops = \textbf{8 XP}
ℹtem \textbf{Totals}: \textbf{32 XP}. With +4 XP, take \emph{Network Builder} (4 XP talent) or add \emph{Minor informant ring} (4 XP).
\end{itemize}

\noindent\textbf{Career Arc:} Early on, the Schemer is just ``the one with
connections.'' Mid game grows their network and enemies. Late game, they might
operate entire webs of spies, informants, merchants, or political factions.

\paragraph{Reminder:}
All builds above assume baseline \emph{Attributes at 1} and \emph{Skills at 0}
before spending. Attribute and Skill advances are cumulative by step (see costs
in \ref{sec:xp−starting−guidelines}).

Remember: Your chosen path should reflect both your character concept and your
preferred play style. There is no single ``correct'' path−−−only what works for
you, your group, and the story you want to tell.

% ============================================================
% LEGACY ENGINE CROSS−REFERENCE
% ============================================================

\subsection*{Planning for Succession (The Legacy Engine)}
If you are playing a long campaign or enjoy the idea of your character's
influence outlasting them, consider investing in a \textbf{Follower} (Cap 3+)
early. That follower can become your \textbf{successor} when your character
retires or falls −− inheriting your faction standing, one unresolved burden, and
a Legacy Bond (see \textbf{Chapter \ref{ch:legacy−engine}: The Legacy Engine}).
A loyal apprentice, a trusted lieutenant, or a family member can carry your
story forward.

% ============================================================
% NARRATIVE−HEAVY BUILDING
% ============================================================

\section{Narrative−Heavy Character Building Options}
łabel{sec:xp−narrative}
∈dex{narrative advancement}

For groups that prefer a strong narrative focus in character building, you can
treat XP as a tool for telling the story of a career, not just a currency.

\paragraph{Story−Driven Milestones}
łabel{sec:xp−milestones}
Instead of tracking XP numerically, the GM can award advancement when
characters reach significant story milestones.  
``You have trained with the master for months−−−you've improved your skill.''

\paragraph{Experience Through Reflection}
łabel{sec:xp−reflection}
Players can spend downtime scenes reflecting on past experiences to earn XP.
A meaningful flashback, confession, or revelation can justify growth without
counting every point.

\paragraph{Collaborative Advancement}
łabel{sec:xp−collaborative}
The group can discuss and agree on advancement choices, ensuring everyone's
growth supports the overall story direction and shared arcs.

\paragraph{Narrative Justification Focus}
łabel{sec:xp−narrative−justification}
When spending XP, players should explain how their character gained this
capability through in−game experiences, creating richer backstory and
continuity. ``I raise Presence because I've been forced to handle the group's
negotiations for weeks.''

\paragraph{Path as Theme}
łabel{sec:xp−path−theme}
Use your chosen path as a \textbf{theme} rather than a strict budget:
\begin{itemize}
  ℹtem A Personal Path character emphasizes physical, mental, or spiritual self−mastery.
  ℹtem A Balanced Path character emphasizes learning how to juggle tools and relationships.
  ℹtem An Influencer Path character emphasizes their growing web of relationships and obligations.
\end{itemize}

In all cases, let your character's \textbf{career} on the page match the career
they are living in the fiction.